% Source: https://artofproblemsolving.com/community/c3h1984706_more_efficient_way_to_approach_this_problem
% Stable identity: aops-topic-1984706-q01
Write down the five $2$'s; they create $6$ gaps (before, between and after). Put $g_i$ ones in gap $i$. A valid number corresponds to exactly one solution of
\[
g_1+g_2+\cdots+g_6=5,\qquad 0\le g_i\le2 .
\]
Without the upper bound there are $\dbinom{10}{5}=252$ solutions. If $g_i\ge3$ for a given $i$, subtract $3$ from it: $\dbinom{7}{5}=21$ solutions, and two gaps cannot both have $3$ or more. By inclusion--exclusion the count is
\[
252-6\cdot21=\boxed{126}.
\]
(Equivalently, it is the coefficient of $x^5$ in $(1+x+x^2)^6$.) The method works for every problem of this type: separate the restricted symbols by the unrestricted ones and count bounded distributions into the gaps.
